% Hình 3.3 — Kiểm chứng bản phát hành cô lập và công bố nguyên tử
\begin{tikzpicture}[node distance=5mm, s/.style={mpbox=37mm, minimum height=19mm}]
  % Hàng 1: chuẩn bị và dựng ứng viên
  \node[s] (p1) {\textbf{1. Phiếu chuẩn bị}\\ ghim snapshot Bronze\\ \texttt{identity} = SHA-256};
  \node[s, right=of p1] (p2) {\textbf{2. Xuất snapshot}\\ file JSON + manifest\\ có checksum, lưu MinIO};
  \node[s, right=of p2] (p3) {\textbf{3. Nạp ứng viên}\\ CSDL \texttt{mp\_release\_*}\\ 8 bảng \texttt{raw\_*}};
  \node[s, right=of p3] (p4) {\textbf{4. dbt build}\\ 9 mô hình\\ 31 bài kiểm thử};
  % Hàng 2: đối soát và công bố (đi ngược từ phải sang trái)
  \node[s, below=8mm of p4] (p5) {\textbf{5. Đối soát checksum}\\ checksum từng bảng\\ oracle CTR/CPC\\ \texttt{proof\_digest}};
  \node[s, left=of p5] (p6) {\textbf{6. Ghi sổ cái}\\ \texttt{certified\_bundles}\\ \texttt{publications}\\ (PostgreSQL)};
  \node[s, left=of p6] (p7) {\textbf{7. Công bố}\\ \texttt{publication\_journal}\\ (ClickHouse)};
  \node[mphead=37mm, minimum height=19mm, left=of p7] (p8) {Ứng dụng đọc\\ \texttt{max(generation)}};
  % Nhánh thất bại
  \node[mpbox=60mm, below=8mm of p6, xshift=-21mm, draw=mpline, fill=mpzone] (fail)
    {\textbf{Bất kỳ bước nào thất bại}\\ ứng viên bị loại bỏ, release đang phục vụ giữ nguyên};

  \draw[mparrow] (p1) -- (p2);
  \draw[mparrow] (p2) -- (p3);
  \draw[mparrow] (p3) -- (p4);
  \draw[mparrow] (p4) -- (p5);
  \draw[mparrow] (p5) -- (p6);
  \draw[mparrow] (p6) -- (p7);
  \draw[mparrow] (p7) -- (p8);
  \draw[mpdash] (p5.south) |- (fail.east);
\end{tikzpicture}
